% Folder 144, batch 02 (archivists' pages 21-40).
% Pass: Opus 5.5 (claude-opus-5-5), 2026-10-03 — first pass, unchecked against the pages by a human.
%
% His own pagination stands on some sheets (6) on p. 22, 8) on p. 24; then
% (1)-5 on pp. 35-39). The batch opens in mid-sentence on p. 21. Page 28 is a
% cover sheet in his hand (its words become the section heading); page 40 is
% blank. Both are skipped.

\documentclass[11pt,a4paper]{article}
\input{../preamble/grothendieck}

\folder{144}
\batch{2}
\pages{21}{40}
\foldertitle{[Suite autour de Teichmüller dont] Teichmülleries, 1981-1982 : notes manuscrites (1981-1983, s.d.), lettres (1981, s.d.).}
\dating{1981-1983}
\watermark{Édition de démonstration}

\begin{document}

\page{21}
\note{la page commence au milieu d'une phrase, ouverte avant le lot.}
« défini sur $K$, pour \ill{} assez grand ». Choisissons \ill{} $K = K$\uncertain{.}, et \uncertain{conjuguons} \ill{} d'exprimer les \uncertain{conjugaisons} \ill{} des types « galoisiens » que l'exprime. Soit $\overline{K}$ la clôture alg. de $K$ dans $k$. \uncertain{Considérons}
\note{les signes $K = K$ de la deuxième ligne sont lus tels quels ; le mot qui les précède n'est pas lu.}

\begin{tikzcd}[column sep=small, row sep=small, nodes={font=\scriptsize}]
\mathfrak{X}_{g,\nu\,\mathbb{Q}} \arrow[d] & \mathfrak{X}_{g,\nu\,K} \arrow[l] \arrow[d] & \mathfrak{X}_{g,\nu\,\overline{K}} \arrow[l, no head] \arrow[d] & \mathfrak{X}_{g,\nu\,k} \arrow[l, no head] \arrow[d] \\
M_{g,\nu\,\mathbb{Q}} \arrow[d] & M_{g,\nu\,K} \arrow[l] \arrow[d] \arrow[u, dashed, bend right=40] & M_{g,\nu\,\overline{K}} \arrow[l, no head] \arrow[d] & M_{g,\nu\,k} \arrow[l, no head] \arrow[d] \\
\mathbb{Q} \arrow[r, no head] & K \arrow[r, no head] \arrow[u, dashed, bend left=40] & \overline{K} \arrow[r, no head] & k
\end{tikzcd}

\note{diagramme numéroté (7). Les deux flèches pointillées partent de $K$ (sections) ; leurs têtes et leurs points d'attache sont lus sur un croquis serré. Les traits horizontaux sans tête sont ceux de la page.}
\marginal{accolade sous $\overline{K} \;—\; k$ : au pt de vue homotopique, cette \uncertain{extension} est de nature triviale}

$(X,x)$ provient d'une \struck{\ill{}} \add{section} (\uncertain{tributaire} \ill{} $x$) de \struck{$\mathfrak{D}$}$\mathfrak{X}_{g,\nu\,K}$ sur $K$, définissant une section (fonction \ill{} $\xi$) de $M_{g,\nu\,K}$ sur $K$, $\{$ les sections $x$, $\xi$ sur $k$, (dont \ill{} \uncertain{se déduisent} par \uncertain{changement} des corps de base).

Posons
\[
(8)\quad
\begin{cases}
\mathcal{M}_{X,x,K} = \pi_1(\mathfrak{X}_{g,\nu\,K}, x) \\
\mathcal{N}_{X,x,K} = \pi_1(M_{g,\nu\,K}, \xi) \\
\Gamma_{\overline{K}|K} = \mathrm{Gal}(\overline{K}|K) = \pi_1(\operatorname{Spec} K, \overline{k})
\end{cases}
\]
\note{dans la première ligne, le premier argument de $\pi_1$ est écrit sur un $M$ biffé, et le point de base $x$ sur un signe biffé ; le numéro (8) est lu avec doute.}

on a cette fois-ci, comme \ill{} (1) \ill{}, $K$ remplaçant $\mathbb{Q}$ :

\page{22}
\note{p. 6) de l'auteur.}

\begin{tikzcd}[column sep=small, row sep=small, nodes={font=\scriptsize}]
 & 1 \arrow[d] & 1 \arrow[d] & & \\
 & \pi_{X,x} \arrow[r, no head, "="] \arrow[d] & \pi_{X,x} \arrow[d] & & \\
1 \arrow[r] & \mathfrak{G}_{X,x} \arrow[r] \arrow[d] & \mathcal{M}_{X,x,K} \arrow[r] \arrow[d] & \Gamma_{\overline{K}|K} \arrow[r] \arrow[d, no head, "\parallel"] & 1 \\
1 \arrow[r] & \mathcal{T}_{X} \arrow[r] \arrow[d] & \mathcal{N}_{X,x,K} \arrow[r] \arrow[d] & \Gamma_{\overline{K}|K} \arrow[r] & 1 \\
 & 1 & 1 & &
\end{tikzcd}

\note{diagramme numéroté (9). La lettre $\mathfrak{G}$ rend une majuscule bouclée qu'il emploie dans tout le lot ; $\mathcal{T}$ rend une majuscule cursive (groupe de Teichmüller) ; l'indice de $\mathcal{T}$ est lu $X$ ou $x$.}

correspondant :
\[
(10)\qquad
\overbrace{\pi_{X,x} \subset \underbrace{\mathfrak{G}_{X,x} \subset \mathcal{M}_{X,x,K}}_{\Gamma_{\overline{K}|K}}}^{\mathcal{N}_{X,K}}
\]
\note{sur la page, une accolade sous $\pi_{X,x} \subset \mathfrak{G}_{X,x}$ porte $\mathcal{T}_X$, une autre sous $\mathfrak{G}_{X,x} \subset \mathcal{M}_{X,x,K}$ porte $\Gamma_{\overline{K}|K}$, et l'accolade du dessus, marquée $\mathcal{N}_{X,K}$, couvre $\mathfrak{G}_{X,x} \subset \mathcal{M}_{X,x,K}$ ; la mise en forme ci-dessus ne rend pas exactement l'emboîtement.}

$\{$\textbf{NB} On n'a pas encore utilisé que $(X,x)$ est défini sur $K$, i.e. que $x$ (provienne \ill{} \add{$\in \mathfrak{X}_{g,\nu}(K)$} \ill{} dans $K$ …). À \uncertain{vrai} dire, la situation (9) \add{— (10)} se déduit de la situation « absolue » (i.e. sur $\mathbb{Q}$) (1), par \struck{la} image inverse des structures d'extension (de $\Gamma_{\widetilde{\mathbb{Q}}/\mathbb{Q}}$ par $\mathfrak{G}_{X,x} \supset \pi_{X,x}$) : \uncertain{l'inclusion} de l'hom. des \ill{}
\[
(11)\qquad \Gamma_{\overline{K}|K} \longrightarrow \Gamma_{\widetilde{\mathbb{Q}}/\mathbb{Q}}
\]

\page{23}
par
\[
(12)\quad
\begin{cases}
\mathcal{M}_{X,x,K} \simeq \mathcal{M}_{X,x} \times_{\Gamma_{\widetilde{\mathbb{Q}}/\mathbb{Q}}} \Gamma_{\overline{K}/K} \\
\mathcal{N}_{X,x,K} \simeq \mathcal{N}_{X,x} \times_{\Gamma_{\widetilde{\mathbb{Q}}/\mathbb{Q}}} \Gamma_{\overline{K}/K}
\end{cases}
\]
\struck{d'où}

\begin{tikzcd}[column sep=small, row sep=small, nodes={font=\scriptsize}]
1 \arrow[r] & \mathfrak{G}_{X,x} \arrow[r] & \mathcal{M}_{X,x} \arrow[r] & \Gamma_{\widetilde{\mathbb{Q}}/\mathbb{Q}} \arrow[r] & 1 \\
1 \arrow[r] & \mathfrak{G}_{X,x} \arrow[r] \arrow[u, no head, "\parallel"] & \mathcal{M}_{X,x,K} \arrow[r] \arrow[u, "\mathrm{cart}"'] & \Gamma_{\overline{K}/K} \arrow[r] \arrow[u] & 1
\end{tikzcd}

\note{diagramme numéroté (13). La lettre $\mathfrak{G}$ de la ligne du haut est tracée sur une autre, biffée.}

et de même

\begin{tikzcd}[column sep=small, row sep=small, nodes={font=\scriptsize}]
1 \arrow[r] & \mathcal{T}_{X,x} \arrow[r] & \mathcal{N}_{X,x} \arrow[r] & \Gamma_{\widetilde{\mathbb{Q}}/\mathbb{Q}} \\
1 \arrow[r] & \mathcal{T}_{X,x} \arrow[r] \arrow[u, no head, "\parallel"] & \mathcal{N}_{X,x,K} \arrow[r] \arrow[u, "\mathrm{cart}"'] & \Gamma_{\overline{K}/K} \arrow[u]
\end{tikzcd}

\note{diagramme numéroté (13'). Le $\mathcal{T}$ du haut est écrit sur une lettre biffée ; le $\mathcal{N}_{X,x,K}$ du bas est écrit sur un $\mathcal{M}$.}

(déduit des précédents en divisant par $\pi_{X,x} \subset \mathfrak{G}_{X,x}$)

Ceci posé, la donnée de la section $x$ de \struck{$M$}$\mathfrak{X}_{g,\nu\,K}$ sur $K$ définissant $x$ sur $k$ (i.e. de la \uncertain{restriction} des corps de base de $k$ à $K$ : $K$ dans $(X,x)$) s'exprime par un \uncertain{scindage} (\add{\ill{}}) de l'ext. $\mathcal{M}_{X,x,K}$,

\page{24}
\note{p. 8) de l'auteur.}
\ill{} donc par un relèvement

\begin{tikzcd}[column sep=small, row sep=small]
\mathcal{M}_{X,x} \arrow[r] & \Gamma_{\widetilde{\mathbb{Q}}/\mathbb{Q}} \\
 & \Gamma_{\overline{K}/K} \arrow[u] \arrow[ul, dashed]
\end{tikzcd}

\note{diagramme numéroté (14).}

et de $=$ pour $\xi$, correspondant à une \uncertain{restriction} du corps de définition de $X$ à $K$, \uncertain{ici} correspondant :

\begin{tikzcd}[column sep=small, row sep=small]
\mathcal{N}_{X,x} \arrow[r, two heads] & \Gamma_{\overline{\mathbb{Q}}/\mathbb{Q}} \\
 & \Gamma_{\overline{K}/K} \arrow[u] \arrow[ul, dashed]
\end{tikzcd}

\note{diagramme numéroté (14').}

\struck{On \uncertain{trouve} donc, puisque}

\marginal{trois traits verticaux dans la marge}
Je pense que ces données permettent de reconstituer $(K,\overline{K},(X,x))$ resp. $(K,\overline{K},X)$ \add{resp. (14')} [au-dessus de (14) :] (14).

J'ai envie d'exprimer la structure \ill{} \uncertain{plus simples}, en notant que (14) \add{resp. (14')} implique
\[
(15)\qquad \Gamma_{\overline{K}|K} \longrightarrow \operatorname{Aut}_{\mathrm{lac}} \pi_{X,x} \qquad \text{resp.}
\]
\[
(15')\qquad \Gamma_{\overline{K}/K} \longrightarrow \operatorname{Autext}_{\mathrm{lac}}(\pi_{X,x})
\]
\note{l'indice de $\operatorname{Aut}$ et de $\operatorname{Autext}$ est un petit mot lu « lac » sans aucune certitude ; il revient p. 25.}

\page{25}
J'ai envie d'abord de montrer comment ces données, \struck{à partir des données près} (d'une \ill{} $\Gamma = \Gamma_{\overline{K}/K}$, \uncertain{ouvert} \add{\ill{}} dans $\Gamma_{\overline{\mathbb{Q}}/\mathbb{Q}} \simeq \Gamma_{\widetilde{\mathbb{Q}}/\mathbb{Q}}$), $\Gamma$ opérant sur un gr. profini : le cas des $\pi_{X,x}$) permet, à \ill{} près, de \uncertain{reconstituer} les \ill{} \struck{Notons que} \add{Il faut de plus} \ill{} des \ill{} \uncertain{des} $\pi_{X,x}$, \struck{l'image de $\Gamma_{\overline{K}/K}$} \add{de} stable par l'action de $\Gamma$, i.e. un \ill{} des discrétisations \ill{} \ill{} en un torseur sous $\mathcal{M}_{g,\nu}' \subset \operatorname{Aut}_{\mathrm{lac}}(\hat{\pi}_{g,\nu})$. Cela donne \ill{} un \ill{} en torseurs sous $\mathcal{M}_{g,\nu}'$ (\uncertain{voire} sous $\mathcal{M}_{g,\nu}$), ou (dans le cas (5')) en torseurs sous $\mathcal{N}_{g,\nu}'$ — qui sont des quotients de $\mathcal{M}_{g,\nu}$ resp. $\mathcal{N}_{g,\nu}$.
\note{page rapide, à demi lisible ; les formules et les mots « torseur », « discrétisations », « stable par l'action de $\Gamma$ » sont sûrs, la prose qui les relie ne l'est pas.}

\page{26}
$K$ corps de t.f. sur $\mathbb{Q}$

$\overline{K}$ clôture algébrique

$D$-structure sur $\overline{K}$ : a) $\overline{\mathbb{Q}} \subset \overline{K}$ ; b) pour \uncertain{toute} courbe alg. $X$ sur $\overline{K}$ de type $g,\nu$, la donnée d'une discrétification de $\pi_1(X)$.

$D$ structure \add{spéciale} : \uncertain{définissable} par plongement $\varphi : \overline{K} \to \mathbb{C}$

\struck{\textbf{Prop.} Soient $\Theta$, $\Theta'$ deux}

Soient $\overline{K}$, $\overline{K}'$ deux clôtures alg. Pour tout $K$-\struck{\ill{}}isom (plus gén., pour tout $K$-isom) $\overline{K} \to \overline{K}'$, \uncertain{par transport de} structure des $D$ structures (resp. spéciales) sur $\overline{K}$ en définit une sur $\overline{K}'$. En particulier, $\Gamma = \operatorname{Aut}(\overline{K}/K)$ opère sur l'ens. des discrétifications (resp. spéciales) sur $\overline{K}$.

\textbf{Prop} Soient $\overline{K}$, $\overline{K}'$ deux clôtures algébriques de $K$, munies de discrét. \emph{spéciales} $\Theta$, $\Theta'$. Alors il existe un $K$-isom $\overline{K} \to \overline{K}'$ compatible avec $\Theta$, $\Theta'$.

\textbf{Cor} Deux discrétifications spéciales sur $\overline{K}$ sont isomorphes par un élément de $\Gamma = \operatorname{Aut}_K(\overline{K})$ (\ill{} \add{\ill{}} d'un sous-gr. d'indice fini \uncertain{dans} $\Gamma$)

\textbf{NB}

\begin{tikzcd}[column sep=small, row sep=small, nodes={font=\scriptsize}]
\operatorname{Isom}_{*}(\overline{K},\overline{K}') \times \operatorname{Disc}(\overline{K}) \arrow[r] & \operatorname{Disc}(\overline{K}') \\
\operatorname{Isom}_{*}(\overline{K},\overline{K}') \times \operatorname{Hom}(\overline{K},\mathbb{C}) \arrow[r] \arrow[u] & \operatorname{Hom}(\overline{K}',\mathbb{C}) \arrow[u]
\end{tikzcd}

\marginal{« ! » et trois traits verticaux dans la marge}
\struck{\ill{}} \add{\ill{}} \uncertain{est commutatif} \struck{\ill{}} si $\overline{K} = \overline{K}'$, \ill{} \ill{} à la limite. $\operatorname{Isom}_K(\overline{K},\overline{K}') = \Gamma$ — \ill{} $\Gamma$ \uncertain{équivariant}.
\note{l'indice $*$ de $\operatorname{Isom}$ est une étoile appuyée, lue aussi $K$ dans la dernière ligne.}

\page{27}
\[
\operatorname{Hom}(\overline{K},\mathbb{C}) \xrightarrow{\ \text{surj}\ } \operatorname{Discspec}(\overline{K})
\]
\struck{\ill{}} compatible avec les op. naturelles de $\Gamma$

\textbf{NB} \struck{\ill{}} $\Gamma$ est transitif$^{*}$, i.e. \struck{droite} … \ill{} transitif : \struck{\ill{}} \add{$\varphi$ \ill{}} (pour que $\varphi$, $\varphi' : \overline{K} \rightrightarrows \mathbb{C}$ soient conjugués par $\Gamma$, il f. et il s. a) $\varphi|K = \varphi'|K$ et b) $\varphi(\overline{K}) = \varphi'(\overline{K})$, conditions qui sont des plus \uncertain{draconiennes} ….

\marginal{en biais, au coin supérieur gauche : « Soit » coché, « alors Soit » ; puis, renvoyé par l'astérisque : $^{*}$\ill{} si $\overline{K} = \overline{\mathbb{Q}}$ i.e. pour $K$ alg. sur $\mathbb{Q}$ ; plus bas : donc \uncertain{explicitons} \ill{} d'injectif}

\emph{Autour} d'une discrétification sur $\overline{K}$ : \ill{} \ill{} \ill{} « le cas » transcendant d'\uncertain{intéressant} pas \struck{\ill{}} des \struck{\ill{}} « le cas » transcendants, sur $K \otimes_{\mathbb{Q}} \mathbb{C}$, \ill{}

Est-il possible que $\operatorname{Hom}(\overline{K},\mathbb{C}) \to \operatorname{Discspec}(\overline{K})$ soit injectif (i.e. bijectif) ? \struck{\ill{} \ill{} \ill{} $\overline{K} \subset \mathbb{C}$ \ill{}} \struck{\ill{}} \ill{} i.e. des $(K, \varphi_K : K \to \mathbb{C})$, \emph{\uncertain{équivalent}} : elles \ill{} \emph{\uncertain{rigidifications}}, et \ill{} des $(K,\overline{K},\Theta)$, $\Theta = \Theta_{\varphi_{\overline{K}}}$ \uncertain{similaire} dans le cas \uncertain{où} $\overline{K}$ alg. sur $\mathbb{Q}$.
\note{passage enchevêtré : des traits courbes relient les morceaux de la phrase et un long trait barre une ligne entière ; l'ordre de lecture retenu est conjectural.}

$\pi_1$, $\pi_{1,\mathrm{ab}}$ \emph{structure de Hodge mixte}
\[
\pi_{1,\mathrm{ab}} \otimes_{\mathbb{Z}} \mathbb{C}
\]
ou encore structure complexe sur $\widetilde{\pi}_{1,\mathrm{ab}} \otimes_{\mathbb{Z}} \mathbb{R}$

\section{Formulaire pour $\pi_{0,3}$ etc. Pt de vue transcendant pour $M_{1,1}$}
\note{titre de sa main, seul inscrit sur un feuillet de couleur (p. 28), qui sert de couverture à ce qui suit.}

\page{29}
\emph{$\mathrm{Gl}(2,\mathbb{Z})$} (Formulaire 2\supplied{e} version) (1982)
\note{la date « (1982) » est d'une autre encre, plus pâle.}

\[
\boxed{\rho = \begin{pmatrix} 0 & 1 \\ -1 & 1 \end{pmatrix}}, \quad
\rho^{-1} = \begin{pmatrix} 1 & -1 \\ 1 & 0 \end{pmatrix}
\qquad \rho^3 = \rho^{-3} = \omega \quad \text{donc } \rho^2 = \omega\rho^{-1}, \quad \rho^6 = 1
\]
\[
\omega = \begin{pmatrix} -1 & 0 \\ 0 & -1 \end{pmatrix} \quad (\text{donc } \omega g = -g) \qquad \omega^2 = 1
\]
\note{une flèche courbe relie $\omega$ à $\rho$ dans la marge gauche.}

\[
\begin{cases}
\boxed{\sigma_\infty = \begin{pmatrix} 0 & 1 \\ -1 & 0 \end{pmatrix}} \qquad
\sigma_\infty^{-1} = \omega\sigma_\infty = \begin{pmatrix} 0 & -1 \\ 1 & 0 \end{pmatrix} \\[4pt]
\sigma_0 = \begin{pmatrix} -1 & 1 \\ -2 & 1 \end{pmatrix} \qquad \rho(\sigma_i) = \sigma_{i+1} \\[4pt]
\sigma_1 = \begin{pmatrix} -1 & 2 \\ -1 & 1 \end{pmatrix}
\end{cases}
\]
\[
\sigma^2 = \sigma^{-2} = \omega, \quad \sigma^{-1} = \omega\sigma \qquad\qquad \sigma_i^2 = \omega, \quad \sigma_i^{-1} = \omega\sigma_i
\]
\note{l'indice de $\sigma_1$ est écrit sur un autre, empâté.}

\[
\begin{cases}
\boxed{\varepsilon_0 = \sigma_\infty^{-1}\rho = \begin{pmatrix} 1 & -1 \\ 0 & 1 \end{pmatrix}} \quad (\text{unipotent \& régulier}) \\[4pt]
\boxed{\varepsilon_1 = \rho\sigma_\infty^{-1} = \begin{pmatrix} 1 & 0 \\ 1 & 1 \end{pmatrix}} \\[4pt]
\varepsilon_\infty = \rho^2\sigma_\infty^{-1}\rho^{-1} = \begin{pmatrix} 2 & -1 \\ 1 & 0 \end{pmatrix}
\end{cases}
\qquad
\begin{cases}
\varrho_0 = \varepsilon_0^2 = \begin{pmatrix} 1 & -2 \\ 0 & 1 \end{pmatrix} \\[4pt]
\varrho_1 = \varepsilon_1^2 = \begin{pmatrix} 1 & 0 \\ 2 & 1 \end{pmatrix} \\[4pt]
\varrho_\infty = \varepsilon_\infty^2 = \begin{pmatrix} 3 & -2 \\ 2 & -1 \end{pmatrix}
\end{cases}
\]
\note{la lettre bouclée qu'il distingue de $\rho$ est rendue ici $\varrho$. Dans la boîte de $\varepsilon_1$, la formule $\rho\sigma_\infty^{-1}$ est écrite en surcharge ; le $\rho$ est empâté.}
\marginal{dans la marge gauche : $N_0 = \begin{pmatrix} 0 & -1 \\ 0 & 0 \end{pmatrix}$, $N_1 = \begin{pmatrix} 0 & 0 \\ 1 & 0 \end{pmatrix}$, $N_\infty = \begin{pmatrix} 1 & -1 \\ 1 & -1 \end{pmatrix}$}

\struck{\ill{}}
\[
\varepsilon_i = \sigma_i^{-1}\rho, \qquad \rho(\varepsilon_i) = \varepsilon_{i+1}, \qquad \rho(\varrho_i) = \varrho_{i+1}
\]
\[
\varrho_\infty\varrho_1\varrho_0 = \omega, \qquad
\varrho'_\infty\varrho'_1\varrho'_0 = 1 \quad \text{où } \varrho'_i = \omega\varrho_i
\]
\[
\begin{cases}
\sigma_i\sigma_{i+1} = \omega\rho\varrho_i = \omega\varrho_{i+1}\rho \\
\rho(\sigma_i) = \sigma_{i+1} = \omega\varrho_{i-1}\varepsilon_{i+1} \\
\sigma_i(\rho) = \rho^{-1}\varrho'_{i-1}
\end{cases}
\]

\[
\begin{cases}
\boxed{\tau_\infty = \begin{pmatrix} -1 & 0 \\ 0 & +1 \end{pmatrix}} \\[4pt]
\tau_0 = \begin{pmatrix} 1 & 0 \\ 2 & -1 \end{pmatrix} \\[4pt]
\tau_1 = \begin{pmatrix} 1 & -2 \\ 0 & -1 \end{pmatrix}
\end{cases}
\qquad
\tau_i^2 = 1 \qquad \rho(\tau_i) = \tau_{i+1}
\]
\[
\tau_i\tau_{i-1} = \underbrace{\omega\varrho_{i+1}}_{\varrho'_{i+1}} \quad \text{i.e.} \quad
\begin{cases}
\tau_1\tau_0 = \varrho'_\infty \\
\tau_\infty\tau_1 = \varrho'_0 \\
\tau_0\tau_\infty = \varrho'_1
\end{cases}
\]
\note{dans l'accolade de droite, un $\omega$ biffé précède chaque $\varrho'$.}

\[
\begin{cases}
\boxed{\tau'_\infty = \sigma_\infty\tau_\infty = \omega\tau_\infty\sigma_\infty = \begin{pmatrix} 0 & 1 \\ 1 & 0 \end{pmatrix}} \\[4pt]
\tau'_0 = \sigma_0\tau_0 = \omega\tau_0\sigma_0 = \begin{pmatrix} 1 & -1 \\ 0 & -1 \end{pmatrix} \\[4pt]
\tau'_1 = \sigma_1\tau_1 = \omega\tau_1\sigma_1 = \begin{pmatrix} -1 & 0 \\ -1 & 1 \end{pmatrix}
\end{cases}
\]
\[
\tau_i'^{\,2} = 1 \qquad \rho(\tau'_i) = \tau'_{i+1}
\]
\[
\tau'_0\tau'_1 = \tau'_1\tau'_\infty = \tau'_\infty\tau'_0 = \omega\rho, \qquad
\tau'_1\tau'_0 = \tau'_\infty\tau'_1 = \tau'_0\tau'_\infty = \omega\rho^{-1}
\]
\[
\tau'_i = \sigma_i\tau_i = \omega\tau_i\sigma_i = \tau_i\sigma_i^{-1}, \qquad
\tau_i = \tau'_i\sigma_i = \omega\sigma_i\tau'_i = \sigma_i^{-1}\tau'_i,
\]
\[
\sigma_i = \tau'_i\tau_i = \omega\tau_i\tau'_i
\]
\note{devant $\tau_i'^{\,2} = 1$, une surcharge biffée. Dans $\tau'_0\tau'_1 = \dots = \omega\rho$, le $\omega$ est écrit sur une autre lettre. La dernière ligne est coupée par le bord du feuillet et lue avec doute.}

\page{30}
\[
\left.
\begin{aligned}
\tau_i(\varepsilon_j) &= \varepsilon_j^{-1} \\
\tau_i(\varrho_j) &= \varrho_j^{-1} \\
\tau_i(\varrho'_j) &= \varrho_j'^{\,-1}
\end{aligned}
\right| \ \text{si } i \neq j
\]
\[
\begin{cases}
\tau_i(\sigma_i) = \sigma_i^{-1} = \omega\sigma_i \qquad \text{[au-dessus :] } \tau'_i(\sigma_i) = \\
\tau'_i(\rho) = \rho^{-1} \\
\tau_i(\rho) = \varrho_{i-1}'^{\,-1}\rho = \sigma_i(\rho^{-1})
\end{cases}
\]
\note{une grande accolade, dans la marge gauche, embrasse les deux groupes de formules ci-dessus.}

\struck{\ill{}}
\[
\begin{aligned}
&\tau_\infty(\varepsilon_0) = \varepsilon_0^{-1} &\qquad& \tau_\infty(\varrho_0) = \varrho_0^{-1} \\
&\tau_\infty(\varepsilon_1) = \varepsilon_1^{-1} && \tau_\infty(\varrho_1) = \varrho_1^{-1} \\
&\tau_\infty(\varepsilon_\infty) = \varrho_0(\varepsilon_\infty^{-1}) && \tau_\infty(\varrho_\infty) = \varrho_0\varrho_1\omega \\
&\phantom{\tau_\infty(\varepsilon_\infty)} = \varrho_1^{-1}(\varepsilon_\infty^{-1}) && \phantom{\tau_\infty(\varrho_\infty)} = \varrho_1^{-1}(\varrho_\infty^{-1}) \\
& && \phantom{\tau_\infty(\varrho_\infty)} = \varrho_0^{-1}(\varrho_\infty^{-1})
\end{aligned}
\]
\note{dans $\varrho_0(\varepsilon_\infty^{-1})$, un exposant $-1$ est ajouté sous $\varrho_0$ ; l'exposant de la dernière ligne et la fin de ligne, empâtée au bord, sont lus avec doute.}

\[
\begin{aligned}
&\sigma_i(\varrho_{i+1}) = \varrho_{i+2} &\qquad& \sigma_i(\varepsilon_{i+1}) = \varepsilon_{i+2} \\
&\sigma_i(\varrho_{i+2}) = \varrho_{i+1} && \sigma_i(\varepsilon_{i+2}) = \varepsilon_{i+1}) \\
&\sigma_i(\varrho_i) = \omega\underbrace{(\varrho_{i+1}\varrho_{i+2})^{-1}}_{\omega\varrho_{i+2}^{-1}\varrho_{i+1}^{-1}} = \varrho_{i+1}(\varrho_i) && \\
& \phantom{\sigma_i(\varrho_i)} = \varrho_{i+2}^{-1}(\varrho_i) &&
\end{aligned}
\]
\note{les indices de cette colonne sont serrés et lus en partie d'après la symétrie des formules.}

— — —

\[
\begin{cases}
\lambda_0 = \varepsilon_0, \ \lambda_1 = \sigma_\infty, \ \lambda_\infty = \rho^{-1} \\
\lambda_\infty\lambda_1\lambda_0 = 1 \qquad \lambda_1^4 = \lambda_\infty^6 = 1
\end{cases}
\qquad
\overbrace{\Gamma_0 = \tau'_\infty}^{= \left(\begin{smallmatrix} 0 & 1 \\ 1 & 0 \end{smallmatrix}\right)}, \
\overbrace{\Gamma_1 = -\tau'_0}^{= \left(\begin{smallmatrix} -1 & 1 \\ 0 & 1 \end{smallmatrix}\right)}, \
\overbrace{\Gamma_\infty = \tau_\infty}^{= \left(\begin{smallmatrix} -1 & 0 \\ 0 & 1 \end{smallmatrix}\right)}
\]
\note{devant $\lambda_0$, une lettre biffée. La lettre $\Gamma_i$ rend une majuscule gamma qu'il trace comme un petit $r$ anguleux ; on pourrait aussi lire $r_i$.}
\[
\Gamma_\infty\Gamma_1 = \lambda_0, \qquad \Gamma_0\Gamma_\infty = \lambda_1, \qquad \Gamma_1\Gamma_0 = \lambda_\infty
\]

Donc \uncertain{on a}

\begin{tikzcd}
\pi^{\tau}_{0,3} \arrow[r, "\varphi"] \arrow[d] & \mathrm{Gl}(2,\mathbb{Z}) \arrow[d] \\
\pi_{0,3} \arrow[r, "\varphi^{0}"] & \mathrm{Sl}(2,\mathbb{Z})
\end{tikzcd}

$\varphi$, $\varphi^{0}$ \uncertain{surj.} \qquad $\varphi(\varrho_i) = \lambda_i$, \quad $\varphi(\tau_i) = \Gamma_i$
\note{les flèches verticales du carré sont tracées comme des crochets ; leur sens n'est pas marqué d'une tête. La lettre sur le $\mathrm{Gl}$ du haut est écrite sur une autre.}

\[
\mathrm{Sl}(2,\mathbb{Z}) = \{\rho, \sigma \mid \rho^3 = \rho^{-3} = \sigma^2 \ (= \sigma^{-2})\}
= \mathbb{Z}/6\mathbb{Z} *_{\mathbb{Z}/2\mathbb{Z}} \mathbb{Z}/4\mathbb{Z}
\]
\[
\begin{aligned}
\mathrm{Gl}(2,\mathbb{Z}) &= \{\rho, \sigma, \tau' \mid \rho^3 = \rho^{-3} = \sigma^2 \ (= \sigma^{-2}),\ \tau'(\rho) = \rho^{-1},\ \tau'(\sigma) = \sigma^{-1}\} \\
&= \{\Gamma_0, \Gamma_1, \Gamma_\infty \mid \Gamma_0^2 = \Gamma_1^2 = \Gamma_\infty^2 = 1,\
(\Gamma_0\Gamma_\infty)^2 = (\Gamma_1\Gamma_0)^3 = (\Gamma_1\Gamma_0)^{-3}\} \\
&= D_6 *_{D_2} D_4
\end{aligned}
\]
\marginal{dans la marge gauche : $\sigma = \sigma_\infty$, $\tau' = \tau'_\infty$}
\note{dans la présentation de $\mathrm{Gl}$, le générateur $\tau'$ est écrit sur un $\tau'_\infty$ empâté ; le premier générateur de la deuxième ligne est écrit $\tau_0$ sur la page.}
\[
\begin{cases}
\text{NB } \tau' = \Gamma_0, \quad \Gamma_0(\rho) = \rho^{-1}, \ \Gamma_0(\sigma) = \sigma^{-1} \\
\Gamma_1 = \rho^{-1}\Gamma_0 = \Gamma_0\rho \\
\Gamma_\infty = \sigma\Gamma_0 = \Gamma_0\sigma^{-1}
\end{cases}
\]
\note{deux biffures empâtées dans cette accolade.}

\emph{Relations avec le formulaire 1981}

Les $\sigma_i$ sont devenus \uncertain{transformés} en $\sigma_i^{-1}$, donc $\varepsilon_0 = \sigma_i^{-1}\rho$ au lieu de $\varepsilon_0 = \sigma_i\rho$. Les autres \uncertain{quantités} pas changées, mais noté $\varrho_i = \varepsilon_i^2$ \struck{\ill{}} au lieu de $h_i$, et $\varrho'_i = \omega\varepsilon_i^2 = \omega\varrho_i$ au lieu de $\lambda_i$ ($\varrho'_\infty\varrho'_1\varrho'_0 = 1$, $\varrho_\infty\varrho_1\varrho_0 = \omega$), où $\omega = -\mathrm{id}$. Les $\lambda'_i$ sont devenus $\lambda_i$.
\note{les deux lettres que la version 1981 employait (« $h_i$ », « $\lambda_i$ ») sont lues avec doute.}

\page{31}
(1982)
\note{date d'une autre encre, au coin supérieur gauche.}
\[
1 \to \mathbb{Z} \xrightarrow{\ k\ } \mathrm{Gl}(2,\mathbb{Z})^{\sim} \to \mathrm{Gl}(2,\mathbb{Z}) \to 1
\]
$k(\mathbb{Z})$ central dans $\mathrm{Sl}(2,\mathbb{Z})^{\sim}$ mais non dans $\mathrm{Gl}(2,\mathbb{Z})^{\sim}$.

$\mathrm{Gl}(2,\mathbb{Z})$ opère sur $\mathbb{Z}$ par l'intermédiaire du dét $u$, par \struck{\ill{}} \add{\uncertain{multiplication}} par dét $u \in \{\pm 1\}$.

\[
k(1) \overset{\mathrm{def}}{=} \ell'_0
\]
\[
\left.
\begin{aligned}
&\tilde{\rho} \text{ au dessus de } \rho, & \tilde{\rho}^{6} &= \ell_0'^{\,-1} \\
&\tilde{\sigma} \text{ — — } \sigma, & \tilde{\sigma}^{4} &= \ell'_0 \\
&\tilde{\omega}_0 \text{ — — } \omega_0 = -1, & \tilde{\omega}^{2} &= \ell'_0
\end{aligned}
\right|
\ \text{conditions sur } \tilde{\rho}, \tilde{\sigma}, \tilde{\omega}_0 \text{ donc :}
\ \left|\ \tilde{\rho}^{-3} = \tilde{\sigma}^{2} = \tilde{\omega}_0 \right.
\]
\marginal{dans un coin encadré à droite : $\mathrm{Gl}(2,\mathbb{Z})$ opère sur $\mathrm{Sl}(2,\mathbb{Z})^{\sim}$ (mais non sur $\mathrm{Gl}(2,\mathbb{Z})^{\sim}$)}

$\tilde{\tau}'$ au dessus de $\tau'$, on aura nécessairement
\[
\begin{cases}
\tilde{\tau}'(\tilde{\rho}) = \tilde{\rho}^{-1}, \quad \tilde{\tau}'(\tilde{\sigma}) = \tilde{\sigma}^{-1}, \quad \tilde{\tau}'^{\,2} = 1 \\
\tilde{\tau}'(\tilde{\omega}_0) = \tilde{\omega}_0^{-1}, \quad \tilde{\tau}'(\ell'_0) = \ell_0'^{\,-1}
\end{cases}
\]
\marginal{à droite, après un trait : \emph{NB} on peut faire \ill{}}

\struck{On \uncertain{posera} \ill{}}

On pose
\[
\tilde{\varepsilon}_0 = \tilde{\sigma}\tilde{\rho}, \qquad
\tilde{\varepsilon}_1 = \tilde{\rho}(\tilde{\varepsilon}_0) = \tilde{\sigma}^{-1}(\varepsilon_0) = \tilde{\sigma}(\varepsilon_0) = \tilde{\rho}\tilde{\sigma}
\]
de sorte qu'on \uncertain{a}
\[
\tilde{\varepsilon}_1\tilde{\varepsilon}_0 = \tilde{\rho}\tilde{\sigma}^2\tilde{\rho} = \tilde{\rho}^{-1} \qquad (\text{car } \tilde{\sigma}^2 = \tilde{\rho}^{-3}), \ \text{donc}
\]
\[
\tilde{\rho} = (\tilde{\varepsilon}_1\tilde{\varepsilon}_0)^{-1} = \tilde{\varepsilon}_0^{-1}\tilde{\varepsilon}_1^{-1}, \qquad
\tilde{\sigma} = \tilde{\rho}^{-1}\tilde{\varepsilon}_1 = \tilde{\varepsilon}_1\tilde{\varepsilon}_0\tilde{\varepsilon}_1 = \tilde{\varepsilon}_0\tilde{\rho}^{-1} = \tilde{\varepsilon}_0\tilde{\varepsilon}_1\tilde{\varepsilon}_0
\]

\[
\mathrm{Sl}(2,\mathbb{Z})^{\sim} \ \left|\
\begin{aligned}
&\tilde{\rho}, \tilde{\sigma} : && \tilde{\rho}^3\tilde{\sigma}^2 = 1 \\
&\tilde{\varepsilon}_0, \tilde{\varepsilon}_1 : && \boxed{\tilde{\varepsilon}_0\tilde{\varepsilon}_1\tilde{\varepsilon}_0 = \tilde{\varepsilon}_1\tilde{\varepsilon}_0\tilde{\varepsilon}_1}
\end{aligned}
\right.
\]
\note{avant la relation encadrée, une première forme, $(\dots)^3 = (\tilde{\varepsilon}_1\tilde{\varepsilon}_0\tilde{\varepsilon}_1)$ ou approchant, est biffée de hachures.}

\[
\mathrm{Gl}(2,\mathbb{Z})^{\sim} \ \left|\
\begin{aligned}
&\tilde{\rho}, \tilde{\sigma}, \tilde{\tau}' = \tilde{\Gamma}_0 && \tilde{\Gamma}_0^{2} = \tilde{\rho}^3\tilde{\sigma}^2 = 1, \ \tilde{\Gamma}_0(\tilde{\rho}) = \tilde{\rho}^{-1}, \ \tilde{\Gamma}_0(\tilde{\sigma}) = \tilde{\sigma}^{-1} \\
&\tilde{\varepsilon}_0, \tilde{\varepsilon}_1, \tilde{\Gamma}_0 && \tilde{\Gamma}_0^{2} = 1, \ \tilde{\varepsilon}_0\tilde{\varepsilon}_1\tilde{\varepsilon}_0 = \tilde{\varepsilon}_1\tilde{\varepsilon}_0\tilde{\varepsilon}_1, \\
& && \tilde{\Gamma}_0(\tilde{\varepsilon}_0) = \tilde{\varepsilon}_1^{-1}, \ \tilde{\Gamma}_0(\tilde{\varepsilon}_1) = \tilde{\varepsilon}_0^{-1} \\
&\tilde{\Gamma}_0, \tilde{\Gamma}_1, \tilde{\Gamma}_\infty && \tilde{\Gamma}_0^2 = \tilde{\Gamma}_1^2 = \tilde{\Gamma}_\infty^2 = (\Gamma_1\Gamma_0)^3(\Gamma_0\Gamma_\infty)^{2} = 1
\end{aligned}
\right.
\]
\note{le « 2 » de $\mathrm{Gl}(2,\mathbb{Z})^{\sim}$ est écrit sur un autre chiffre ; devant la première relation, une surcharge biffée. L'exposant de $(\Gamma_0\Gamma_\infty)$ est lu $2$ (peut-être $-2$).}

\[
\begin{aligned}
&\tilde{\varepsilon}_i \ (i \in \{0,1,\infty\}) \quad \rho(\tilde{\varepsilon}_i) = \tilde{\varepsilon}_{i+1} &\quad& \tilde{\ell}_i = \tilde{\varepsilon}_i^{\,2}, \ \rho(\tilde{\ell}_i) = \tilde{\ell}_{i+1} \\
&\tilde{\sigma}_i \ (i \in \{0,1,\infty\}) \quad \rho(\tilde{\sigma}_i) = \tilde{\sigma}_{i+1} = \tilde{\ell}_{i-1}\tilde{\varepsilon}_{i+1} && \tilde{\sigma}_i^{\,2} = \tilde{\rho}^{-3} = \tilde{\omega}_0, \quad \tilde{\sigma}_i^{\,4} = \tilde{\rho}^{-6} = \ell'_0
\end{aligned}
\]
\[
\tilde{\varepsilon}_i = \tilde{\sigma}_{i-1}\tilde{\rho} = \tilde{\rho}\tilde{\sigma}_{i+1}
\]
\[
\tilde{\sigma}_i\tilde{\sigma}_{i+1} = \tilde{\omega}_0\tilde{\rho}\tilde{\ell}_i = \tilde{\omega}_0\tilde{\ell}_{i+1}\tilde{\rho}
\qquad
\tilde{\sigma}_i(\tilde{\rho}) = \tilde{\rho}^{-1}\tilde{\lambda}_{i-1} = \tilde{\omega}_0^{-1}\tilde{\rho}^{-1}\tilde{\ell}_{i-1}
\]
\[
\tilde{\ell}_\infty\tilde{\ell}_1\tilde{\ell}_0 = \tilde{\omega}_0 \quad \text{i.e.} \quad
\tilde{\lambda}_\infty\tilde{\lambda}_1\tilde{\lambda}_0 = \tilde{\omega}_0^{-2} = \ell_0'^{\,-1}
\quad (\text{où } \tilde{\lambda}_i = \tilde{\omega}_0^{-1}\tilde{\ell}_i)
\]
\[
\text{i.e.} \quad \tilde{\lambda}_0^{-1}\tilde{\lambda}_1^{-1}\tilde{\lambda}_\infty^{-1} = \ell'_0
\]
\note{les tildes de cette page sont posés de façon irrégulière ; on les restitue là où la page les porte nettement. Les lettres $\ell$, $\lambda$ des dernières lignes sont mal distinctes, et une lettre biffée sépare $\tilde{\rho}^{-1}$ de $\tilde{\lambda}_{i-1}$.}

\page{32}
\[
\tilde{\rho}^{3}(\tilde{\tau}'_\infty) = \text{\struck{\ill{}}}\,, \qquad \tilde{\omega}_0^{-1}(\tilde{\tau}'_\infty) = \ell_0'^{\,-1} \cdot \tilde{\tau}'_\infty
\]
La suite bi-infinie des \struck{\ill{}} $\tilde{\rho}^{n}(\tilde{\tau}'_\infty)$ ($n \in \mathbb{Z}$) prennent \uncertain{injectivement} l'une des relevés de $\{\tau'_\infty, \tau'_0, \tau'_1\}$. \uncertain{Idem} pour $(\tau_\infty, \tau_0, \tau_1)$.

\[
\begin{cases} \tilde{\Gamma}_0(\tilde{u}) = \tilde{v} \\ \tilde{\Gamma}_0(\tilde{v}) = \tilde{u} \end{cases}
\qquad
\begin{cases} \tilde{\Gamma}_\infty(\tilde{u}) = \tilde{u}^{-1} \\ \tilde{\Gamma}_\infty(\tilde{v}) = \tilde{u}(\tilde{v}) \end{cases}
\qquad
\begin{cases} \tilde{\Gamma}_1(\tilde{u}) = \tilde{u}^{-1} \\ \tilde{\Gamma}_1(\tilde{v}) = \tilde{u}\tilde{v} \end{cases}
\]
\[
\begin{cases} \tilde{\rho}(\tilde{u}) = \tilde{v}^{-1} \\ \tilde{\rho}(\tilde{v}) = \tilde{v}\tilde{u} \end{cases}
\qquad
\begin{cases} \tilde{\varepsilon}_0(\tilde{u}) = \tilde{u} \\ \tilde{\varepsilon}_0(\tilde{v}) = \tilde{v}\tilde{u}^{-1} \end{cases}
\qquad
\begin{cases} \tilde{\sigma}(\tilde{u}) = \tilde{u}(\tilde{v}) \\ \tilde{\sigma}(\tilde{v}) = \tilde{u}^{-1} \end{cases}
\qquad
\begin{cases} \tilde{\varepsilon}_1(\tilde{u}) = \tilde{u}\tilde{v} \\ \tilde{\varepsilon}_1(\tilde{v}) = \tilde{v} \end{cases}
\]
\note{le reste de la page est blanc. Le sens de « $\tilde{u}(\tilde{v})$ » (action, ou produit avec parenthèse) n'est pas fixé par la page.}

\page{33}
\[
\boxed{\mathfrak{S}_{0,3} \xrightarrow{\ \sim\ } \mathrm{Gl}(2,\hat{\mathbb{Z}})'} \qquad
\boxed{\mathfrak{S}^{+}_{0,3} \simeq \mathrm{Sl}(2,\hat{\mathbb{Z}})'} \qquad \text{formulaire (1981)}
\]
\note{la lettre $\mathfrak{S}$ rend une majuscule gothique ; le chapeau de $\hat{\mathbb{Z}}$ est clair dans le premier encadré, moins dans le second. Les astérisques en marge de certaines lignes sont de lui.}

\[
\begin{aligned}
&{*}\ \rho = \begin{pmatrix} 0 & 1 \\ -1 & 1 \end{pmatrix} \quad \rho^{-1} = \begin{pmatrix} 1 & -1 \\ 1 & 0 \end{pmatrix}
&\quad& \rho^3 = -1 \ \text{i.e.}\ \rho^2 = -\rho^{-1} \qquad \rho^6 = 1 \\
&{*}\ \sigma_\infty = \begin{pmatrix} 0 & -1 \\ 1 & 0 \end{pmatrix} && \sigma_i^2 = -1 \ \text{i.e.}\ \sigma_i^{-1} = -\sigma_i \qquad \sigma_i^4 = 1 \\
&\phantom{*}\ \sigma_0 = \begin{pmatrix} 1 & -1 \\ 2 & -1 \end{pmatrix} \\
&\phantom{*}\ \sigma_1 = \begin{pmatrix} 1 & -2 \\ 1 & -1 \end{pmatrix} && \rho\sigma_i\rho^{-1} = \sigma_{i+1} \qquad \sigma_i\sigma_{i+1} = \rho\ell_i = -\ell_{i+1}\rho
\end{aligned}
\]
\[
\begin{cases}
\rho(\sigma_i) = \sigma_{i+1} = \ell_{i-1}\varepsilon_{i+1} \\
\sigma_i(\rho) = \rho^{-1}\lambda_{i-1}
\end{cases}
\]
\note{dans $\rho = \dots$, une flèche est biffée avant le signe $=$ ; dans $\sigma_i^2 = -1$, le signe est écrit sur un autre. L'accolade ci-dessus est d'une encre plus foncée, ajoutée.}

\[
\left|
\begin{aligned}
&{*}\ \varepsilon_0 = \sigma_\infty\rho = \begin{pmatrix} 1 & -1 \\ 0 & 1 \end{pmatrix} = 1 + N_0 &\quad& \ell_0 = \varepsilon_0^2 = \begin{pmatrix} 1 & -2 \\ 0 & 1 \end{pmatrix} \\
&{*}\ \varepsilon_1 = \rho\sigma_\infty = \begin{pmatrix} 1 & 0 \\ 1 & 1 \end{pmatrix} = 1 + N_1 && \ell_1 = \varepsilon_1^2 = \begin{pmatrix} 1 & 0 \\ 2 & 1 \end{pmatrix} \\
&\phantom{*}\ \varepsilon_\infty = \rho^2\sigma_\infty\rho^{-1} = \begin{pmatrix} 2 & -1 \\ 1 & 0 \end{pmatrix} = 1 + N_\infty && \ell_\infty = \varepsilon_\infty^2 = \begin{pmatrix} 3 & -2 \\ 2 & -1 \end{pmatrix}
\end{aligned}
\right.
\]
\marginal{dans la marge gauche, d'une autre encre : $N_0 = \begin{pmatrix} 0 & -1 \\ 0 & 0 \end{pmatrix}$, $N_1 = \begin{pmatrix} 0 & 0 \\ 1 & 0 \end{pmatrix}$, $N_\infty = \begin{pmatrix} 1 & -1 \\ 1 & -1 \end{pmatrix}$}
\[
\rho\varepsilon_i\rho^{-1} = \varepsilon_{i+1}, \qquad \rho\ell_i\rho^{-1} = \ell_{i+1}
\]
\[
\ell_\infty\ell_1\ell_0 = -1 \quad \text{i.e.} \quad \lambda_\infty\lambda_1\lambda_0 = 1 \quad (\lambda_i = -\ell_i), \qquad
\varepsilon_i = \sigma_{i-1}\rho = \rho\sigma_{i+1}
\]

\[
\begin{aligned}
&{*}\ \tau_\infty = \begin{pmatrix} -1 & 0 \\ 0 & +1 \end{pmatrix} &\quad& \tau_i^2 = 1 \\
&\phantom{*}\ \tau_0 = \begin{pmatrix} -1 & 0 \\ -2 & -1 \end{pmatrix} && \tau_i\tau_{i+1} = -\ell_{i+2} \ \text{i.e.} \ 
\begin{cases} \tau_1\tau_0 = -\ell_\infty \\ \tau_\infty\tau_1 = -\ell_0 \\ \tau_0\tau_\infty = -\ell_1 \end{cases} \\
&\phantom{*}\ \tau_1 = \begin{pmatrix} -1 & -2 \\ 0 & -1 \end{pmatrix} && \rho\tau_i\rho^{-1} = \tau_{i+1}
\end{aligned}
\]
\note{les matrices de $\tau_0$ et $\tau_1$ sont lues telles qu'écrites ($\tau_0$ : $-1, 0 / -2, -1$ ; $\tau_1$ : $-1, -2 / 0, -1$) ; elles ne sont pas d'ordre $2$, et ne s'accordent pas avec la version de 1982 (p. 29). Dans $\tau_i\tau_{i+1} = -\ell_{i+2}$ l'indice est lu d'après les trois relations explicites ; une première colonne des mêmes relations est biffée de hachures.}

\[
\begin{aligned}
&{*}\ \tau'_\infty = \text{\struck{\ill{}}}\,\sigma_\infty\tau_\infty = +\tau_\infty\sigma_\infty = \begin{pmatrix} 0 & 1 \\ 1 & 0 \end{pmatrix} \\
&{*}\ \tau'_0 = -\sigma_0\tau_0 = +\tau_0\sigma_0 = \begin{pmatrix} 1 & -1 \\ 0 & -1 \end{pmatrix} \\
&\phantom{*}\ \tau'_1 = -\sigma_1\tau_1 = +\tau_1\sigma_1 = \begin{pmatrix} -1 & 0 \\ -1 & 1 \end{pmatrix}
\end{aligned}
\qquad
\begin{cases}
\tau_i'^{\,2} = 1 \\
\rho\tau'_i\rho^{-1} = \tau'_{i+1} \\
\tau'_0\tau'_1 = \tau'_1\tau'_\infty = \tau'_\infty\tau'_0 = -\rho \\
\tau'_i(\rho) = \rho^{-1} \quad \tau'_i(\sigma_i) = \tau_i(\sigma_i) = \sigma_i^{-1} = -\sigma_i
\end{cases}
\]
\struck{\ill{}}, $\tau_i(\varepsilon_j) = \varepsilon_j^{-1}$ si $i \neq j$
\[
\left.
\begin{aligned}
\tau_i(\sigma_i) &= \sigma_i^{-1} = -\sigma_i \\
\tau_i(\rho) &= -\ell_{i-1}^{-1}\rho = \sigma_i(\rho^{-1})
\end{aligned}
\right|
\quad
\begin{cases}
\tau'_i = -\sigma_i\tau_i = +\tau_i\sigma_i \\
\tau_i = \tau'_i\sigma_i = +\sigma_i\tau'_i \\
\sigma_i = -\tau'_i\tau_i = +\tau_i\tau'_i
\end{cases}
\]
\note{dans la deuxième ligne de droite, des signes sont surchargés et biffés ; la lecture retenue est conjecturale.}

$\mathrm{Sl}(2,\mathbb{Z})$ présenté par (NB $\sigma = \sigma_\infty$)
\[
\boxed{\begin{cases} \sigma^4 = \rho^6 = 1 \\ \sigma^2 = \rho^3 \end{cases}}
\]
$\mathrm{Gl}(2,\mathbb{Z})$ présenté par
\[
\boxed{\begin{cases} \Gamma_0^2 = \Gamma_1^2 = \Gamma_\infty^2 = 1 \\ (\Gamma_0\Gamma_\infty)^4 = (\Gamma_1\Gamma_0)^6 = 1 \\ (\Gamma_0\Gamma_\infty)^2 = (\Gamma_1\Gamma_0)^3 \end{cases}}
\qquad \simeq \ D_6 \amalg_{D_2} D_4
\]
\note{les indices des $\Gamma$ dans l'encadré sont repassés à l'encre foncée sur d'autres.}

\[
\begin{cases}
\Gamma_0(\sigma_\infty) = \sigma_\infty^{-1} \\
\Gamma_0(\rho) = \rho^{-1} \\
\Gamma_1 = \rho^{-1}\Gamma_0 = \Gamma_0\rho \\
\Gamma_\infty = \sigma\Gamma_0 = \Gamma_0\sigma^{-1}
\end{cases}
\]

\textbf{NB}
\[
\begin{cases}
\ell'_0 = \varepsilon_0, \ \ell'_1 = (-\sigma_\infty), \ \ell'_\infty = \rho^{-1} \\
\ell'_\infty\ell'_1\ell'_0 = 1
\end{cases}
\qquad
\Gamma_0 = \tau'_\infty, \ \Gamma_1 = -\tau'_0, \ \Gamma_\infty = \tau_\infty
\]
\[
\begin{cases}
\tau_\infty\tau'_0 = \varepsilon_0 \\
\tau'_\infty\tau_\infty = -\sigma_\infty \\
(-\tau'_0)\tau'_\infty = \rho^{-1}
\end{cases}
\ \text{i.e.} \
\begin{cases}
\Gamma_\infty\Gamma_1 = \ell'_0 \\
\Gamma_0\Gamma_\infty = \ell'_1 \\
\Gamma_1\Gamma_0 = \ell'_\infty
\end{cases}
\]
\marginal{d'une encre pâle, avec une flèche vers ces lignes : « attention aux signes ! »}
\note{$(-\sigma_\infty)$ est écrit au-dessus d'une valeur biffée ; dans la troisième ligne de gauche, la valeur biffée avant $\rho^{-1}$ est illisible. Un mot biffé suit l'accolade.}

\[
\begin{cases}
\tau_i(\varepsilon_j) = \varepsilon_j^{-1} \ \text{si } i \neq j \\
\tau'_i(\varepsilon_{i+1}) = \varepsilon_{i-1}^{-1} \\
\sigma_i(\varepsilon_{i+1}) = \varepsilon_{i-1}
\end{cases}
\qquad
\begin{cases}
\sigma_i(\ell_{i+1}) = \ell_{i-1} \\
\sigma_i(\ell_{i-1}) = \ell_{i+1}
\end{cases}
\]
\[
\sigma_i(\ell_i) = -(\ell_{i+1}\ell_{i-1})^{-1} \ (\text{attention aux signes !})
\]
\[
\begin{cases}
\tau'_i(\ell_{i-1}) = \ell_{i+1}^{-1} \\
\tau'_i(\ell_{i+1}) = \ell_{i-1}^{-1}
\end{cases}
\ \tau'_i(\ell_i) = \ell_i^{-1}
\qquad
\begin{cases}
\tau_i(\ell_{i-1}) = \ell_{i-1}^{-1} \\
\tau_i(\ell_{i+1}) = \ell_{i+1}^{-1}
\end{cases}
\ \tau_i(\ell_i) = -(\ell_{i-1}^{-1}\ell_{i+1}^{-1})^{-1} = -\ell_{i+1}\ell_{i-1}
\]
\marginal{en bas à droite, encre pâle : « attention aux signes ! »}
\note{les deux dernières lignes sont d'une encre pâle, très serrées ; les indices sont lus en partie par symétrie.}

\page{34}
\[
[\ell_0, \ell_1] = [\lambda_0, \lambda_1] = (\sigma_\infty(\rho)\rho^{-1})^3 \qquad \sigma_\infty(\rho)\rho^{-1} = \ell_0\rho = \rho\ell_\infty
\]
\[
[\sigma_\infty, \ell_\infty] = [\sigma_\infty, \lambda_\infty] = (\sigma_\infty(\rho^{-1})\rho)^3 = \mathrm{int}(\sigma)(\rho^{-1}) \cdot \underbrace{[\ell_0, \ell_1]}_{[\lambda_0, \lambda_1]}
\]
donc
\[
[\lambda_0, \lambda_1] = 1 \iff [\sigma_\infty, \ell_\infty] = 1 \iff (\underbrace{\sigma_\infty(\rho)\rho^{-1}}_{\ell_0\rho = \rho\ell_\infty})^3 = 1
\]
\note{le reste de la page est blanc. La lettre lue $\rho$ dans « $\sigma_\infty(\rho)$ » est tracée comme un $\varphi$ ; la lecture suit la formule de la p. 33 ($\sigma_i(\rho) = \rho^{-1}\lambda_{i-1}$).}

\page{35}
(1) $\varepsilon_0, \varepsilon_1 \in H$ (\add{gén. du} groupe $H$) tels que
\[
\boxed{\varepsilon_0\varepsilon_1\varepsilon_0 = \varepsilon_1\varepsilon_0\varepsilon_1} \quad (*)
\]
\note{après $(*)$, une formule biffée, dont on lit la fin $\dots\varepsilon_1^{-1}\varepsilon_0^{-1}\varepsilon_1$.}

On pose
\[
\begin{cases}
\sigma = \varepsilon_0\varepsilon_1\varepsilon_0 \ (= \varepsilon_1\varepsilon_0\varepsilon_1) \\
\rho'_0 = \sigma\varepsilon_0 = \dots = (\varepsilon_1\varepsilon_0)^2 \\
\rho'_1 = \sigma\varepsilon_1 = \varepsilon_0\sigma = (\varepsilon_0\varepsilon_1)^2 \\
\rho_0 = \sigma^{-1}\varepsilon_0 = \varepsilon_1\sigma^{-1} \ (= \varepsilon_0^{-1}\varepsilon_1^{-1} = (\varepsilon_1\varepsilon_0)^{-1}) \\
\rho_1 = \varepsilon_0\sigma^{-1} = \sigma^{-1}\varepsilon_1 \ (= \varepsilon_1^{-1}\varepsilon_0^{-1} = (\varepsilon_0\varepsilon_1)^{-1})
\end{cases}
\]
\note{dans les lignes de $\rho'_0$ et $\rho'_1$, des égalités intermédiaires sont biffées de hachures ; on garde ce qui reste lisible. Le premier membre de $\varepsilon_0^{-1}\varepsilon_1^{-1}$ est suivi d'un signe empâté. La première lettre, $\rho_0$, est tracée en boucle, comme le $\varrho$ de la p. 29 ; on la rend $\rho$ ici, la page n'opposant plus deux lettres. Dans le crochet qui suit, la relation $\rho_0 = (\varepsilon_1\varepsilon_0)^{-1}$ et la définition $\rho_0 = \sigma^{-1}\varepsilon_0$ ne s'accordent que si $\varepsilon_1\varepsilon_0\varepsilon_1 = \sigma$, ce qui est $(*)$.}

[donc $\sigma(\varepsilon_0) = \varepsilon_1$, $\sigma(\varepsilon_1) = \varepsilon_0$, $\sigma^2 \in \operatorname{Centre}(H)$, $\operatorname{int}(\sigma) = \operatorname{int}(\sigma^{-1})$]
\marginal{à droite, un petit schéma : $\rho_0, \rho'_0$ \quad $\rho_1, \rho'_1$ ; au-dessous $\sigma$ ; au-dessous $\varepsilon_0$ \quad $\varepsilon_1$}

Alors \struck{$\sigma = \varepsilon_0\rho_0$}

\[
\text{(a)} \quad
\begin{cases}
\boxed{\begin{aligned} &\sigma(\varepsilon_0) = \varepsilon_1, \quad \sigma(\varepsilon_1) = \varepsilon_0 \\ &\sigma(\rho_0) = \rho_1, \quad \sigma(\rho_1) = \rho_0 \end{aligned}} \\
\sigma(\rho'_0) = \rho'_1, \quad \sigma(\rho'_1) = \rho'_1
\end{cases}
\]
\marginal{en biais devant l'accolade : « conj. par $\sigma$ »}
\note{la dernière égalité porte bien $\sigma(\rho'_1) = \rho'_1$ sur la page ; on attendrait $\rho'_0$.}

(\emph{NB} \struck{\ill{}} \add{la première} des \uncertain{cinq} relations \uncertain{équivaut} à $(*)$, et elles impliquent [\add{la seconde} \struck{\ill{}}] $\sigma^2$ commute : $\varepsilon_0$, $\varepsilon_1$, donc au groupe engendré \uncertain{donc à} $H$ \uncertain{engendré par} $\varepsilon_0, \varepsilon_1$) $\forall h \in H$, on a aussi $\sigma(h) = \sigma^{-1}(h)$
\note{la reconstitution de cette parenthèse, à deux insertions croisées, est conjecturale.}

\[
\text{(a')} \quad
\begin{cases}
\rho_0(\varepsilon_0) = \rho'_0(\varepsilon_0) = \varepsilon_1 \\
\rho_1(\varepsilon_1) = \rho'_1(\varepsilon_1) = \varepsilon_0 \\[4pt]
\varepsilon_0(\rho_0) = \rho_1, \quad \varepsilon_0(\rho'_0) = \rho'_1 \\
\varepsilon_1(\rho_1) = \rho_0, \quad \varepsilon_1(\rho'_1) = \rho'_0
\end{cases}
\]
\marginal{en biais devant l'accolade : « conj. par $\rho_i, \rho'_i$ » et « conj. par $\varepsilon_i$ »}

\[
\text{(b)} \quad
\begin{aligned}
&\rho_0 = \sigma^{-1}\varepsilon_0 = \varepsilon_1\sigma^{-1}, \quad \rho_1 = \sigma^{-1}\varepsilon_1 = \varepsilon_0\sigma^{-1} \\
&\rho'_0 = \sigma\varepsilon_0 = \varepsilon_1\sigma, \quad \rho'_1 = \sigma\varepsilon_1 = \varepsilon_0\sigma
\end{aligned}
\ \left|\ \text{$\rho_i, \rho'_i$ en termes de $\sigma$, $\varepsilon$.}\right.
\]

\struck{(c) \ill{} $\rho'_0 = \rho_0^{-2}$, $\rho'_1 = \rho_1^{-2}$ \ill{} $\rho_0'^{\,-1}\rho_0 = \rho_0^3$ \ill{} : $\varepsilon_0$ \ill{} $\rho_1'^{\,-1}\rho_1 = \rho_1^3$ \ill{} $\varepsilon_1$)}
\[
\text{(c)} \quad \rho_0^{-3} = \rho_1^{-3} = \sigma^{2} \ \Big(\overset{\mathrm{def}}{=} \omega \in Z(H)\Big)
\]
\[
\text{(d)} \quad \rho'_0 = \rho_0^{-2} = \omega\rho_0, \quad \rho'_1 = \rho_1^{-2} = \omega\rho_1
\qquad [\text{donc } \operatorname{int}(\rho_i) = \operatorname{int}(\rho'_i) = \operatorname{int}(\rho_i)^{-2}]
\]
\[
\text{d'où} \quad \rho_0 = \omega^{-1}\rho'_0 \ (= \rho'_0\omega^{-1}), \quad \rho_1 = \omega^{-1}\rho'_1 = \rho'_1\omega^{-1}
\]
\note{dans (c), l'exposant de $\sigma$ est un $2$ appuyé sur un autre chiffre, et le $\omega$ est écrit sur une lettre biffée. Dans (d), un signe est empâté avant $\rho_0^{-2}$, et une amorce biffée suit la ligne.}

\emph{NB} Il vaut mieux au début ne pas introduire $\rho'_0, \rho'_1$, qu'on \uncertain{a} : ce \uncertain{qui} \uncertain{fera} des données un formulaire complémentaire \uncertain{très} \uncertain{encombrant}

(Juin 1983)

\page{36}
\note{feuillet numéroté (2) par l'auteur, à la suite de la p. 35 (numérotée (1)).}

(e) \quad [$\varepsilon_0, \varepsilon_1$]
\[
\sigma = \varepsilon_0\varepsilon_1\varepsilon_0 = \varepsilon_1\varepsilon_0\varepsilon_1, \quad
\rho_0 = (\varepsilon_1\varepsilon_0)^{-1}, \ \rho_1 = (\varepsilon_0\varepsilon_1)^{-1}, \quad
\rho'_0 = (\varepsilon_1\varepsilon_0)^2, \ \rho'_1 = (\varepsilon_0\varepsilon_1)^2
\]
relation fondamentale
\[
\boxed{\varepsilon_0\varepsilon_1\varepsilon_0 = \varepsilon_1\varepsilon_0\varepsilon_1 \quad \text{ou} \quad \varepsilon_0\varepsilon_1\varepsilon_0\varepsilon_1^{-1}\varepsilon_0^{-1}\varepsilon_1^{-1} = 1}
\]

(f$_0$) \quad [$\sigma, \varepsilon_0$]
\[
\varepsilon_1 = \sigma(\varepsilon_0) = \sigma\varepsilon_0\sigma^{-1}, \quad
\boxed{\begin{aligned} \rho_0 &= \sigma^{-1}\varepsilon_0 \\ \rho_1 &= \varepsilon_0\sigma^{-1} \end{aligned}}, \
\begin{aligned} \rho'_0 &= \sigma\varepsilon_0 \\ \rho'_1 &= \varepsilon_0\sigma \end{aligned}
\qquad
\boxed{\begin{aligned} &[\sigma^2, \varepsilon_0] = 1 \\ &(\sigma^{-1}\varepsilon_0)^3 = \sigma^{-2} \end{aligned}}
\]
\note{l'exposant $-2$ du dernier $\sigma$ est écrit sur un autre chiffre.}

Symétriques pour expressions en $\sigma, \varepsilon_1$ (relations f$_1$)

(g$_0$) \quad [$\sigma, \rho_0$]
\[
\varepsilon_0 = \sigma\rho_0, \ \varepsilon_1 = \rho_0\sigma, \quad
\rho'_0 = \rho_0^{-2}, \quad
\rho_1 = \sigma(\rho_0) = \sigma\rho_0\sigma^{-1},
\]
\[
\rho'_1 = \sigma(\rho_0^{-2}) = \sigma\rho_0^{-2}\sigma^{-1}
\qquad
\boxed{\rho_0^{-3} = \underbrace{\sigma^{2}}_{\omega}}
\]
\struck{\ill{}} élément \struck{\ill{}} \uncertain{nécessairement} central
\note{l'accolade marquée $\omega$ est sous $\rho_0^{-3}$ sur la page ; un mot biffé suit $\rho_0^{-2}$.}

Symétriques pour expressions en $\sigma/\rho_1$ (relations g$_1$)

(g$'_0$) \quad [$\sigma, \rho'_0$]
\[
\varepsilon_0 = \sigma^{-1}\rho'_0, \ \varepsilon_1 = \rho'_0\sigma^{-1}, \quad
\rho_0 = \sigma^{-2}\rho'_0, \quad
\rho_1 = \sigma^{-1}\rho'_0\sigma^{-1} =, \quad
\rho'_1 = \sigma\rho'_0\sigma^{-1}
\qquad
\boxed{\begin{aligned} &[\sigma^2, \rho'_0] = 1 \ \text{et} \\ &\rho_0'^{\,3} = \sigma^4 \end{aligned}}
\]

Symétriques pour expressions en $\sigma, \rho'_1$ (relations g$'_1$)

(h$_0$) \quad [$\varepsilon_0, \rho_0$]
\[
\varepsilon_1 = \rho_0(\varepsilon_0) = \rho_0\varepsilon_0\rho_0^{-1}, \quad
\sigma = \varepsilon_0\rho_0^{-1}, \quad
\rho'_0 = \rho_0^{-2},
\]
\[
\rho_1 = \varepsilon_0(\rho_0) = \varepsilon_0\rho_0\varepsilon_0^{-1}, \quad
\rho'_1 = \varepsilon_0(\rho_0^{-2}) = \varepsilon_0\rho_0^{-2}\varepsilon_0^{-1}
\]
\[
\boxed{\rho_0\varepsilon_0\rho_0^{-1}\varepsilon_0\rho_0 = 1}
\]
\note{l'encadré est doublé d'un cadre en pointillés, comme provisoire. Les crochets [$\dots$] rendent les paires de générateurs qu'il inscrit dans la marge gauche, en regard de chaque cas.}

Symétriques pour expressions en $\varepsilon_1, \rho_1$ (relations h$_1$) \quad (à suivre !)

\page{37}
\note{feuillet numéroté 3 par l'auteur. Une longue flèche, dans la marge, remonte du cas (h$'_0$) vers le cas (i$_0$), qui ouvre la page.}

(i$_0$) \quad [$\varepsilon_0, \rho'_0$]
\[
\varepsilon_1 = \rho'_0(\varepsilon_0) = \rho'_0\varepsilon_0\rho_0'^{\,-1}, \quad
\sigma = \rho'_0\varepsilon_0^{-1}, \quad
\rho_0 = \varepsilon_0\rho_0'^{\,-1}\varepsilon_0, \quad
\rho_1 = \varepsilon_0^2\rho_0'^{\,-1}, \quad
\rho'_1 = \varepsilon_0\rho'_0\varepsilon_0^{-1}
\]
\[
\boxed{\begin{aligned}
&\rho'_0 = (\rho'_0(\varepsilon_0)\varepsilon_0)^2 \\
&\quad \text{i.e.} \ \rho'_0 = (\rho'_0\varepsilon_0\rho_0'^{\,-1}\varepsilon_0)^2 \\
&\text{et} \ [\rho_0'^{\,3}, \varepsilon_0] = 1
\end{aligned}}
\]
\note{dans le crochet de la dernière ligne, un premier signe et un groupe après $\rho_0'^{\,3}$ sont biffés de hachures. L'indice du cas, cerclé, est écrit sur une lettre empâtée.}

\emph{NB} On peut se borner à la seule équation \uncertain{abracadabrante} :
\[
\underbrace{\rho'_0\varepsilon_0\rho_0'^{\,-1}\varepsilon_0\rho'_0}\varepsilon_0\underbrace{\rho_0'^{\,-1}\varepsilon_0^{-1}\rho'_0\varepsilon_0^{-1}\rho_0'^{\,-1}}\varepsilon_0^{-1} = 1
\]
i.e.
\[
[\underbrace{\rho'_0\varepsilon_0\rho_0'^{\,-1}\varepsilon_0\rho'_0}, \varepsilon_0] = 1
\]
qui implique \struck{\ill{}} $\rho'_0 = (\rho'_0(\varepsilon_0)\varepsilon_0)^2$ donc $(\rho'_0(\varepsilon_0)\varepsilon_0)^3$
\note{une flèche marquée « $=$ » va de l'accolade du crochet vers $(\rho'_0(\varepsilon_0)\varepsilon_0)^3$ ; le premier symbole de l'équation, devant $\rho'_0$, est biffé. L'exposant de $\rho'_0$ dans « qui implique » est écrit sur un autre.}

Sym. pour expressions en $\varepsilon_1, \rho'_1$ (relations (j$_1$))

(h$'_0$) \quad [$\varepsilon_0, \rho_1$]
\[
\varepsilon_1 = \rho_1^{-1}(\varepsilon_0) = \rho_1^{-1}\varepsilon_0\rho_1, \quad
\sigma = \rho_1^{-1}\varepsilon_0, \quad
\rho_0 = \varepsilon_0^{-1}(\rho_1) = \varepsilon_0^{-1}\rho_1\varepsilon_0,
\]
\[
\rho'_0 = \text{\struck{\ill{}}} = \varepsilon_0^{-1}(\rho_1^{-2}) = \varepsilon_0^{-1}\rho_1^{-2}\varepsilon_0, \quad
\rho'_1 = \rho_1^{-2}
\qquad
\boxed{\rho_1\varepsilon_0\rho_1^{-1}\varepsilon_0\rho_1 = 1}
\]
\marginal{en biais dans la marge gauche : (Le cas (h$'_0$) se déduit de (h$_0$) par conjugaison par $\varepsilon_0$) ; \emph{NB} Il vaut mieux prendre $\varepsilon_1, \rho_0$}
\note{l'étiquette du cas, cerclée, est écrite à côté d'une autre biffée ; dans la première formule, l'indice de $\varepsilon$ est corrigé (1 en 0 ?). L'encadré est doublé de pointillés.}

Symétriques pour expressions en $(\varepsilon_1, \rho_0)$ (relations h$'_1$)

(j$'_0$) \quad [$\varepsilon_0, \rho'_1$]
\[
\varepsilon_1 = \rho_1'^{\,-1}(\varepsilon_0) = \rho_1'^{\,-1}\varepsilon_0\rho'_1, \quad
\sigma = \varepsilon_0^{-1}\rho'_1, \quad
\rho_1 = \varepsilon_0\rho_1'^{\,-1}\varepsilon_0,
\]
\[
\rho'_0 = \varepsilon_0^{-1}(\rho'_1) = \varepsilon_0^{-1}\rho'_1\varepsilon_0, \quad
\rho_0 = \rho_1'^{\,-1}\varepsilon_0^2
\]
\[
\boxed{\begin{aligned}
&\rho'_1 = (\rho'_1(\varepsilon_0)\varepsilon_0)^2 \\
&\quad \text{i.e.} \ \rho'_1 = (\rho'_1\varepsilon_0\rho_1'^{\,-1}\varepsilon_0)^2 \\
&\text{et} \ [\rho_1'^{\,3}, \varepsilon_0] = 1
\end{aligned}}
\]

Symétriques pour expressions en $(\varepsilon_1, \rho'_0)$ (relations \struck{\ill{}} i$'_0$)
\note{l'étiquette de ce cas est lue « j$'_0$ » ; la page annonce pourtant, au bas, les relations « i$'_0$ » pour la paire symétrique, et le cas du haut est désigné « i$_0$ » dans le cercle mais « j$_1$ » dans son renvoi. On transcrit les étiquettes telles qu'elles sont lues.}

\page{38}
\note{feuillet numéroté 4 par l'auteur.}

\emph{NB} Dans cette présentation, symétrie \uncertain{évidente} par deux aut. involutifs de $H$
\[
\begin{aligned}
&1^\circ)\ g \mapsto \check{g} \ (\text{antiaut.}) &\ \Big|\ & \varepsilon_0 \mapsto \varepsilon_0^{-1}, \ \varepsilon_1 \mapsto \varepsilon_1^{-1}, \quad \sigma \mapsto \sigma^{-1}, \\
& & & \rho_0 \mapsto \rho_1^{-1}, \ \rho_1 \mapsto \rho_0^{-1}, \quad \rho'_0 \mapsto \rho_1'^{\,-1}, \ \rho'_1 \mapsto \rho_0'^{\,-1} \\
&2^\circ)\ g \mapsto g^{*} = \operatorname{int}(\sigma) &\ \Big|\ & \varepsilon_0 \mapsto \varepsilon_1, \ \varepsilon_1 \mapsto \varepsilon_0, \quad \sigma \mapsto \sigma, \\
& & & \rho_0 \mapsto \rho_1, \ \rho_1 \mapsto \rho_0, \quad \rho'_0 \mapsto \rho'_1, \ \rho'_1 \mapsto \rho'_0
\end{aligned}
\]
\note{« antiaut. » est souligné deux fois sur la page.}

Sa symétrie \uncertain{centrée} sur $\sigma$ ou $(\varepsilon_0, \varepsilon_1)$, symétrie binaire ($\sigma^2 \in \operatorname{Cent}(H)$), \uncertain{devient} si on centre la symétrie sur $\rho \overset{\mathrm{def}}{=} \rho_0$, \uncertain{on trouve} une symétrie ternaire ($\rho^3 \in \operatorname{Cent}(H)$), \uncertain{avec}
\[
\varepsilon_0, \quad \varepsilon_1 = \rho(\varepsilon_0), \quad \varepsilon_\infty = \rho(\varepsilon_1) = \rho^2(\varepsilon_0)
\]
\uncertain{éléments} circulairement permutés par $\operatorname{int}(\rho)$,

\begin{tikzcd}[column sep=small, row sep=large]
 & \varepsilon_\infty \arrow[dl, no head, "\sigma_1"'] & \\
\varepsilon_0 \arrow[rr, "\sigma_\infty"'] & & \varepsilon_1 \arrow[ul, "\sigma_0"']
\end{tikzcd}

\[
\begin{aligned}
\sigma_1 &= \varepsilon_\infty\varepsilon_0\varepsilon_\infty = \varepsilon_0\varepsilon_\infty\varepsilon_0, \qquad
\sigma_0 = \varepsilon_1\varepsilon_\infty\varepsilon_1 = \varepsilon_\infty\varepsilon_1\varepsilon_\infty, \\
\sigma_\infty &= \varepsilon_0\varepsilon_1\varepsilon_0 = \varepsilon_1\varepsilon_0\varepsilon_1 \ (= \sigma)
\end{aligned}
\]
\[
\rho = (\varepsilon_1\varepsilon_0)^{-1} = (\varepsilon_\infty\varepsilon_1)^{-1} = (\varepsilon_0\varepsilon_\infty)^{-1}
\]
\note{les formules des $\sigma_i$ sont écrites le long des côtés du triangle ; sur la page, le côté $\varepsilon_0\varepsilon_1$ porte une tête vers $\varepsilon_1$, le côté droit une tête vers $\varepsilon_\infty$, le côté gauche aucune. Sous $\rho = \dots$, un ensemble de formules est biffé d'un gribouillis.}

\emph{NB} Mais les éléments
\[
\rho_1 = (\varepsilon_0\varepsilon_1)^{-1} = \sigma_\infty(\rho), \quad
(\varepsilon_1\varepsilon_\infty)^{-1} = \sigma_0(\rho) = \rho(\rho_1), \quad
(\varepsilon_\infty\varepsilon_0)^{-1} = \sigma_1(\rho) = \rho^2(\rho_1)
\]
sont distincts

Les éléments essentiels du point de vue engendrement de $ST_{11}$ \struck{\ill{}} semblent être finalement
\[
\varepsilon_0, \ \varepsilon_1, \ \sigma, \ \rho \ (= \rho_0)
\qquad
\left|\
\begin{aligned}
&(\varepsilon_0, \varepsilon_1) \\
&(\varepsilon_0, \sigma) \leftrightarrow (\varepsilon_1, \sigma) \\
&(\varepsilon_0, \rho) \leftrightarrow (\varepsilon_1, \rho) \\
&(\sigma, \rho)
\end{aligned}
\right.
\]
plus l'élément $\omega = \rho^{-3} = \sigma^2$ (qu'ils engendrent)
\marginal{en biais à droite des couples : les \uncertain{quatre} \uncertain{systèmes} de gén. principaux de $ST_{11}$}
\note{la lecture « $ST_{11}$ » est sûre pour les lettres ; ce que le sigle désigne n'est pas expliqué sur ces pages. Dans le dernier mot de la marge, un mot est biffé.}

\page{39}
\note{feuillet numéroté 5 par l'auteur ; en haut à droite, une date entre parenthèses lue (1983), le premier chiffre empâté.}

\[
\begin{aligned}
&(1)\ (\varepsilon_0, \varepsilon_1), &&\boxed{\varepsilon_0\varepsilon_1\varepsilon_0 = \varepsilon_1\varepsilon_0\varepsilon_1}
&&\begin{cases} \sigma = \varepsilon_0\varepsilon_1\varepsilon_0 = \varepsilon_1\varepsilon_0\varepsilon_1 \\ \rho = (\varepsilon_1\varepsilon_0)^{-1} \end{cases} \\
&(2)\ (\varepsilon_0, \sigma), &&\boxed{[\sigma^2, \varepsilon_0] = 1, \ (\sigma^{-1}\varepsilon_0)^3 = \sigma^{-2}}
&&\begin{cases} \varepsilon_1 = \sigma(\varepsilon_0) = \sigma\varepsilon_0\sigma^{-1} \\ \rho = \sigma^{-1}\varepsilon_0 \end{cases} \\
&(3)\ (\varepsilon_0, \rho) &&\boxed{\rho\varepsilon_0\rho^{-1}\varepsilon_0\rho = 1}
&&\begin{cases} \varepsilon_1 = \rho(\varepsilon_0) = \rho\varepsilon_0\rho^{-1} \\ \sigma = \varepsilon_0\rho^{-1} \end{cases} \\
&(4)\ (\sigma, \rho) &&\boxed{\rho^{-3} = \sigma^2}
&&\begin{cases} \varepsilon_0 = \sigma\rho \\ \varepsilon_1 = \rho\sigma \end{cases}
\end{aligned}
\]
\marginal{à droite de (2) : [sym. $\varepsilon_0 = \sigma(\varepsilon_1) = \sigma\varepsilon_1\sigma^{-1}$, $\rho = \varepsilon_1\sigma^{-1}$] ; à droite de (3) : [sym. $\varepsilon_0 = \rho^{-1}(\varepsilon_1) = \rho^{-1}\varepsilon_1\rho$, $\sigma = \rho^{-1}\varepsilon_1$]}
\note{les encadrés de (1), (3), (4) sont doubles. Dans (3), $\sigma = \varepsilon_0\rho^{-1}$ : l'indice de $\rho$ est écrit sur une autre lettre.}

\note{tableau récapitulatif, encadré, séparé du haut de la page par un double trait :}
\[
\left[\varepsilon_0, \varepsilon_1, \sigma, \rho\right]
\quad
\begin{array}{|l|l|}
\hline
\begin{aligned}
&\varepsilon_0 = \sigma\rho, \quad \sigma = \varepsilon_0\varepsilon_1\varepsilon_0 = \varepsilon_1\varepsilon_0\varepsilon_1 \\
&\varepsilon_1 = \rho\sigma, \quad \rho = (\varepsilon_1\varepsilon_0)^{-1} = \varepsilon_0^{-1}\varepsilon_1^{-1} \\
&\sigma(\varepsilon_0) = \rho(\varepsilon_0) = \varepsilon_1 \\
&\sigma(\varepsilon_1) = \rho^{-1}(\varepsilon_1) = \varepsilon_0
\end{aligned}
&
\begin{aligned}
&\varepsilon_0\varepsilon_1\varepsilon_0 = \varepsilon_1\varepsilon_0\varepsilon_1 \\
&\rho^{-3} = \sigma^2 \ (\overset{\mathrm{def}}{=} \omega) \ \text{central} \\
&\rho\varepsilon_0\rho^{-1}\varepsilon_0\rho = 1
\end{aligned}
\\ \hline
\end{array}
\]
\note{dans la marge du tableau, après $\rho$, un mot biffé. Dans $\sigma(\varepsilon_1) = \rho^{-1}(\varepsilon_1)$, l'indice du second $\varepsilon$ est écrit sur un autre.}

Pour un formulaire plus complet, il faudrait
a) Introduire aussi $\varepsilon_\infty$, \struck{\ill{}} les $\ell_i = \varepsilon_i^2$ ($i \in \{0, 1, \infty\}$)

b) Introduire aussi les $\tau_i$ ($i \in \{0, 1, \infty\}$) pour $ST_{11}^{\pm}$

c) Donner (pour la relation suppl. $\omega^2 = 1$) la représentation naturelle dans $\mathrm{Gl}(2, \mathbb{Z})$
\note{l'exposant de $ST_{11}$ dans b) est lu $\pm$ ; il pourrait être $\mathrm{I}$.}


\end{document}
